\subsection{Forced affine profiles over a three-by-two minimal band}

\begin{proposition}[A common rational affine origin]
\label{C:act-common-affine-origin-cups}
Let $M$ have a completely simple minimal ideal with row set $I$
and two columns.  Its minimal maximal group may be nontrivial.
Suppose that the row maps of all actors with nonconstant column
map fix a common point
\[
 p\in\mathbb QI,\qquad \sum_{i\in I}p_i=1.
\]
Then all rational degree-two products of every finite act cone
$B\mathcal C(M,S)$ vanish.  The point $p$ need not be a row or
have nonnegative coordinates.
\end{proposition}
\begin{proof}
Use the rational constant-action quotient of
Lemma~\ref{C:constant-action-quotient-H2} and the compression of
Theorem~\ref{C:constant-minimal-act-compression}.
This gives a band $I\times\bar S$ even when the original
minimal maximal group is nontrivial.  As proved in that
compression theorem, the constant value of an old minimal-ideal
element depends only on its column: left multiplication varies
the row and group coordinate arbitrarily while preserving that
value.  The two old columns therefore give the same equivariant
map $j\mapsto s_j$ as in
Theorem~\ref{C:act-common-fixed-row-cups}.
Use its honest Karoubi transfer and affine cocycles over
$\mathbb Q$.
Retain its old base row $i_0$, base column $s_0$, covectors $b_f$,
and augmentation operators $A_f$.  Put $t=p-[i_0]$.
For an old monoid arrow replace its row translation by
\[
 a'_m=a_m+(A_m-1)t=\sigma_m p-p.
\]
Put $a'_f=0$ on cross arrows and on the sink identity.  These
assignments are typed: the monoid object uses affine origin $p$
and the sink uses origin $i_0$.  A cross arrow has constant row
map, so its augmentation operator is zero.  Thus on every
composable pair the same affine identity
$a'_{fg}=a'_f+A_fa'_g$ holds.

The cochain $u'_L(f,g)=b^L_fa'_g$ represents the original
degree-two class.  Indeed, on normalized composable pairs its
difference from $u_L$ is
\[
 b^L_f(A_g-1)t=-\delta h(f,g),\qquad h(f)=b^L_ft,
\]
where $h$ is zero on identities.  This equation includes the
cross-arrow/old-arrow pairs; it follows from
$b^L_{fg}=b^L_g+b^L_fA_g$.
The column character $\chi$ now satisfies
$\chi(m)a'_m=0$: for $\chi(m)\ne0$ this is the fixed-point
hypothesis, and otherwise it is automatic.
Define $\alpha'_L(m)=\ell_L(a'_m)$ and use the same normalized
primitive
\[
 P(f,g,h)=u'_L(f,g)\beta(g)\alpha'_D(h)
\]
on nondegenerate chains, extended by zero on degenerate chains.
The calculation in Theorem~\ref{C:act-common-fixed-row-cups}
has defect
$\chi(h)\beta(g)\alpha'_D(k)b^L_fA_ga'_h=0$.
It therefore proves $\delta P=u'_L\smile u'_D$ over
$\mathbb Q$.  The degree-two quotient surjection proves the
assertion for the original act.
\end{proof}

\begin{lemma}[Four necessary outside profiles]
\label{C:three-row-forced-profiles}
Suppose $M$ has a minimal ideal with three rows and two columns and
some finite right act cone has a nonzero rational degree-two
product.  After relabelling rows, $M$ has actors with each of
the following four distinct row/column profiles:
\[
\begin{array}{c|c|c}
 &\sigma&\tau\\ \hline
 p_0&(0,0,2)&(0,1)\\
 p_1&(1,1,2)&(0,1)\\
 r_0&(2,2,0)&(1,0)\\
 r_1&(2,2,1)&(1,0).
\end{array}
\]
Every actor satisfies the balance equation for row values
$(0,0,1)$.  There are exactly $24$ possible row/column
transformation pairs satisfying that equation.  In particular,
the four displayed profiles and the identity require at least
five elements outside the minimal band.
\end{lemma}
\begin{proof}
First $H^2(BM;\mathbb Q)\ne0$.  Otherwise every degree-two
class of the act cone lifts to cohomology relative to $BM$.
The relative quotient is a suspension of the action space,
whose positive-degree products vanish.  Relative cup naturality
would make every degree-two product of the cone zero.

Collapse its minimal maximal group using
Lemma~\ref{lem:B-collapse-minimal-group}; this preserves rational
cohomology and all row/column actions.  The balanced-tensor
description of $H^2$ of the resulting band monoid supplies
a nonzero covector on the row augmentation module with
\[
 \ell A_m=\chi(m)\ell,\qquad
 \chi(m)=\tau_m(1)-\tau_m(0)\in\{0,1,-1\}.
\]
Write the three row values as $v_i$, up to common translation.
If they are all distinct, an actor with $\chi=1$ fixes all three
rows: an affine translation preserving a finite set is the
identity.  An actor with $\chi=-1$, if present, is the unique
reflection of that three-point set and fixes its middle row.
The preceding proposition applies, a contradiction.
Thus there are two levels.  Relabel and rescale them as
$(v_0,v_1,v_2)=(0,0,1)$.

An actor with $\chi=1$ fixes row $2$ and preserves $\{0,1\}$.
An actor with $\chi=-1$ has row map either $(2,2,0)$ or
$(2,2,1)$, and its column map is the swap.  If neither occurs,
row $2$ is a common fixed row.  If only the first occurs,
every $\chi=1$ actor also fixes row $0$: composing its row map
on the left with the first displayed flip would otherwise
produce the second flip.  Then
$( [0]+[2])/2$ is a common rational affine fixed point of every
actor with nonconstant columns.  The other single-flip case
has fixed point $( [1]+[2])/2$.  Each possibility contradicts
Proposition~\ref{C:act-common-affine-origin-cups}.
Both flips are therefore necessary.  Their squares and mixed
products force $p_0,p_1$.

For completeness, the profile count is literal.  There are four
row maps for the identity column map, two for the swap, and
nine for each constant column map.  In the latter case all row
images must lie in one level, giving $2^3+1=9$ choices.
Hence the count is $4+2+9+9=24$.
Minimal-ideal actors have constant column maps, whereas the four forced
profiles and the global identity do not; the five are distinct.
\end{proof}

\begin{lemma}[Projective row modules obstruct every act]
\label{C:row-pd-one-act-cups}
Let $M$ be a finite monoid with minimal ideal $K$, row set $I$,
contracted rational algebra $A=\mathbb QM/\mathbb QK$, and row
augmentation module $V=\widetilde{\mathbb QI}$.
If $\operatorname{pd}_A V\leq1$, then every finite right
$M$-act cone has zero rational degree-two products.
No bound on the act or its homology is required.
The assertion also permits a nontrivial minimal maximal group.
\end{lemma}
\begin{proof}
Pass to the constant-action quotient $\bar S$.
Lemma~\ref{C:constant-action-quotient-H2} gives a surjection on
rational degree-two cohomology in the pullback direction.
Compress by Theorem~\ref{C:constant-minimal-act-compression}.
The resulting monoid $N$ has minimal band $I\times\bar S$.
Its outside elements and their products modulo the minimal
ideal are literally those of $M$, so its contracted algebra
is $A$ and its row module is $V$.
Writing $\bar W=\widetilde{\mathbb Q\bar S}$, the band
dimension-shift formula gives
\[
 H_4(BN;\mathbb Q)
   =\operatorname{Tor}^{A}_2(\bar W,V)=0.
\]
All degree-two products on $BN$ are zero, and cup naturality
and the degree-two surjection give the result.  If $K=M$,
then $M$ is a group, $V=0$, and the rational group-algebra
averaging argument gives the same conclusion directly.
\end{proof}

\begin{proposition}[Five outside elements, with no state bound]
\label{C:three-row-five-outside-all-acts}
If a finite monoid has a minimal ideal with three rows and two
columns and at most five elements outside it, every finite act cone has zero
rational degree-two products.
\end{proposition}
\begin{proof}
If the necessary profiles of
Lemma~\ref{C:three-row-forced-profiles} are absent, that lemma
already proves the assertion.  Otherwise there are exactly
five outside elements: the identity and the four profiles
$p_0,p_1,r_0,r_1$.  Their products cannot land in the band,
because their column maps are nonconstant, and there is only
one element of each profile.  Their literal multiplication is
\[
 (i,\epsilon)(j,\delta)=(i,\epsilon+\delta),
 \qquad i,j\in\{0,1\},\quad\epsilon,\delta\in C_2,
\]
with $(i,0)=p_i$ and $(i,1)=r_i$.
Thus $A$ is the unitalization of $\mathbb Q[L_2\times C_2]$.
Put
\[
 e_\pm=\frac{p_0\pm r_0}{2},\qquad f=1-p_0,\qquad
 n_\pm=\frac{p_1-p_0\pm(r_1-r_0)}2.
\]
The three idempotents $e_+,e_-,f$ are orthogonal and sum to
one.  The radical has basis $n_+,n_-$, with
$n_\pm\in fAe_\pm$ and radical square zero.  This is the
path algebra of a three-vertex fork, hence hereditary.
More directly the row module itself is $Ae_-$: in the row
basis $[1]-[0],[2]-[0]$, the element $e_-$ projects onto
$[2]-[0]$ and $n_-$ sends that vector to $-[1]+[0]$.
The resulting two-dimensional equivariant map $Ae_-\to V$
is an isomorphism.  Lemma~\ref{C:row-pd-one-act-cups}
finishes the proof.
\end{proof}

\begin{lemma}[The complete forced-profile row certificates]
\label{C:forced-profile-row-certificates}
Among ordinary isomorphism classes of monoids with minimal
band $3\times2$ and all four forced profiles, the following
table is complete.  ``Simple periodic kernel'' means an exact
sequence $0\to T\to P\to V\to0$ with $P$ projective and a
one-dimensional module $T$ admitting
$0\to T\to Q\to T\to0$ with $Q$ projective.
\[
\begin{array}{c|r|r|r|r|r}
 |M|&\text{raw tables}&\text{classes}&
  V\text{ projective}&\operatorname{pd}V\leq1\text{, extra}&
  \text{simple periodic kernel}\\ \hline
 11&1&1&1&0&0\\
 12&27&10&10&0&0\\
 13&786&152&148&4&0\\
 14&32188&2980&2908&41&31
\end{array}
\]
Here raw tables use sorted outside profiles before taking
ordinary row/column/element isomorphisms; opposition is not
used to identify different orientations.
\end{lemma}
\begin{proof}
We give the finite classification and the exact module
certificates, so the scope does not depend on a published census.
Fix identity label $0$ and the six row-major band labels.
For order $n$ select the four required profiles and a multiset
of $n-11$ further profiles from the $24$ possibilities above.
The numbers of multisets for $n=11,12,13,14$ are
$1,24,300,2600$.
The identity and all products involving the band are forced.
For each ordered pair of other elements, try every label whose
profile is their composite, including a band label whenever
allowed.  Associativity propagates known products and rejects
only an incompatible equality.  A finished table is checked
on every triple.  Every actual monoid under consideration is
obtained after relabelling its outside elements, proving
coverage of this search.

The code \path{a_forced_swap_targets.cpp} implements this
enumeration.  Its fast version first sorts the transformed
outside profiles under each row/column permutation and then
tries all permutations within equal-profile fibers.  This is
a complete ordinary-isomorphism invariant; it need not be the
unrestricted lexicographic minimum.
The separate program \path{b_forced_profile_raw.cpp} repeats
the raw enumeration with reverse-cell branching and no producer
canonicalizer.  The separate
\path{b_verify_forced_profile_coverage.py} expands every retained
class under all row, column, and equal-profile permutations.
These literal orbits are disjoint and cover the independently
generated raw tables exactly.  It requires final metadata,
exact counts, and complete line lengths before comparison.
The resulting certificate is
\path{b_forced_profile_coverage_verified.json}.

For the module columns, work over the literal contracted basis
$A=\mathbb QM/\mathbb QK$.  The stored row action is the exact
two-dimensional augmentation action.  Each cover is a left
ideal $A\epsilon$, with
$\epsilon=(e-g)/2$ for a literal local pair
$e^2=e$, $g^2=e$, $eg=ge=g$.
The certificate stores the full basis of this ideal, the
equivariant surjective cover to $V$, and its embedded kernel.
For a one-dimensional kernel with character $\theta$, a
projective section is exhibited by an element $z\in A$ with
\[
 az=\theta(a)z\quad\text{on every basis element},\qquad
 \theta(z)=1.
\]
These equations imply $z^2=z$ and identify the kernel with
the projective left ideal $Az$.
The files \path{a_forced_swap_rational_rows12.json},
\path{a_forced_swap_rational_rows13.json}, and
\path{a_forced_swap_rational_rows14.json} store every cover.
The independent Fraction-elimination verifier
\path{d_verify_forced_swap_row_projects.py} imports no producer
or Flint routines.  It checks the literal monoids, ideals,
contracted products, idempotents, full left ideals, surjectivity,
equivariance, kernels, and all displayed projective sections.

There are exactly $31$ remaining one-dimensional kernels at
order $14$.  For each, the file
\path{a_forced_swap_simple_tails14.json} supplies an idempotent
$\eta$, a basis of the two-dimensional ideal $A\eta$, its
surjection to the kernel character, and a one-dimensional
kernel with exactly that same character.  Thus these are
literal exact sequences $0\to T\to A\eta\to T\to0$,
not a periodicity inference from ranks.  The separate Fraction
verifier \path{code/d_verify_forced_swap_simple_tails.py}
reconstructs every full ideal, cover, and kernel and checks the
character equality without importing producer routines or
trusting a recurrence status.  Its final output is
\path{code/outputs/d_forced_swap_simple_tails_verified.json}.
\end{proof}

\begin{theorem}[Every three-by-two act cone through order fourteen]
\label{C:three-row-fourteen-act-cones}
Let $M$ have minimal band $3\times2$.
If $|M|\leq13$, all rational degree-two products of every
finite right act cone vanish, without any homology bound.
If $|M|\leq14$ and the act cone has bounded rational homology,
its rational degree-two products also vanish.
In particular no act cone over such a monoid realizes
$\mathbb CP^2$, at any number of states.
\end{theorem}
\begin{proof}
If the four forced profiles are absent, use
Lemma~\ref{C:three-row-forced-profiles}.
For at most five outside elements use
Proposition~\ref{C:three-row-five-outside-all-acts}.
The row certificates at orders $12$ and $13$, and all but
$31$ cases at order $14$, have projective dimension at most
one, so Lemma~\ref{C:row-pd-one-act-cups} applies.

In one of the remaining $31$ cases, the two certified exact
sequences imply, for every finite-dimensional right module $W$,
\[
 \operatorname{Tor}^{A}_n(W,V)
 \cong\operatorname{Tor}^{A}_{n-1}(W,T)
 \cong\operatorname{Tor}^{A}_1(W,T)
 \qquad(n\geq2).
\]
If the act has a state fixed by the minimal band, its
degree-two products already vanish by the band-grid relative
argument of Lemma~\ref{C:four-state-band-dichotomy}.
Otherwise Lemma~\ref{C:two-column-projective-orbit-kernel}
identifies the original and constant-action-quotient homology
in every degree at least three.  Compress the latter to $N$.
It has the same contracted algebra $A$ and row module $V$,
so the band dimension shift and the displayed recurrence give
\[
 H_q(BN;\mathbb Q)\cong H_4(BN;\mathbb Q)
 \qquad(q\geq4).
\]
Bounded rational homology forces this common group to be zero.
It follows that the original cone has zero fourth rational
homology and hence zero degree-two products, as required.
\end{proof}

\begin{corollary}[The whole three-by-two branch below twenty-two arrows]
\label{C:three-row-act-cones21}
No act cone $\mathcal C(M,S)$ with minimal band $3\times2$
and $|M|+|S|+1\leq21$ realizes $\mathbb CP^2$.
\end{corollary}
\begin{proof}
The preceding theorem handles $|M|\leq14$ at every state
count.  Otherwise $|S|\leq5$.  For at most two states use
the finite-line act argument; three states require at least
$22$ arrows by Theorem~\ref{C:three-state-cone-minimum22}.
For four states, the band dichotomy gives either a fixed state
and zero products, or constant action and a compressed
$3\times4$ monoid; the latter is excluded by
Theorem~\ref{B:three-by-four-targets}.
For five states apply
Theorem~\ref{C:five-state-three-by-two-cones}.
\end{proof}

\begin{corollary}[Nontrivial minimal groups below twenty-two arrows]
\label{C:nontrivial-minimal-group-act-cones21}
No finite act cone with a nontrivial minimal maximal group and
at most $21$ arrows realizes $\mathbb CP^2$.
\end{corollary}
\begin{proof}
At most two minimal rows are excluded by
Theorem~\ref{C:two-row-act-cups}.  One minimal column makes
$BM$ rationally acyclic, so the relative suspension argument
excludes that case.  At most three states are already excluded
below $22$ arrows.  Thus a remaining example has at least three
rows, two columns, four states, and at most $16$ monoid elements.
If its minimal maximal group has order $h\geq2$, its minimal
ideal has $rch$ elements and is proper.  The only possibility
under this bound is $r=3,c=2,h=2$: every other choice costs at
least $16$ minimal-ideal elements and one identity outside it.
The monoid consequently has at most four outside elements,
contrary to the necessary five of
Lemma~\ref{C:three-row-forced-profiles}.
\end{proof}
